\documentclass[11pt]{article}
\usepackage[margin=1in]{geometry}
\usepackage{amsmath,amssymb,amsthm}
\usepackage{microtype}
\usepackage{xcolor}
\usepackage{listings}
\usepackage{hyperref}
\hypersetup{colorlinks=true, linkcolor=blue, urlcolor=blue, citecolor=blue}

\lstset{
  basicstyle=\ttfamily\small,
  breaklines=true,
  numberstyle=\tiny\color{gray},
  numbers=left,
  frame=single,
  language={},
  showstringspaces=false,
}

\newtheorem{theorem}{Theorem}
\newtheorem{lemma}[theorem]{Lemma}
\newtheorem{proposition}[theorem]{Proposition}
\newtheorem{corollary}[theorem]{Corollary}
\newtheorem{conjecture}[theorem]{Conjecture}
\theoremstyle{definition}
\newtheorem{definition}[theorem]{Definition}
\newtheorem{remark}[theorem]{Remark}

\title{Claude Code as the Agent Loop:\\ Per-Step Verification Rate of LLM-Driven Lean 4 Formalization in Extremal Combinatorics}
\author{Matthew Loftus\\Cedar Loop LLC\\\texttt{matthew.a.loftus@gmail.com}}
\date{\today}

\begin{document}

\maketitle

\begin{abstract}
We present a calibrated pipeline for LLM-driven proof formalization in Lean 4 +
Mathlib4, using Claude Code itself as the agent loop with no external orchestration.
The pipeline is calibrated on Loftus (2026) ``Exact combinatorial results for random
2-orders'' (QG Paper G), a recent paper containing 15+ exact identities for the
uniform distribution on pairs of permutations of $\{0, 1, \ldots, N-1\}$.
We formalize 5 of 10 selected identities (C1: ordered comparable pairs;
C2: expected ordering fraction; C4 at $k=2$: 2-element antichains; C5: poset
dimension equals 2 with probability $1 - 1/N!$; C8: expected interval interior
size) over 6 sessions of pair-programming with Claude Code. The aggregate per-step
verification rate is $\approx 63\%$ across $30+$ sub-lemmas, well above the
$20$--$30\%$ floor reported in contemporaneous LLM-for-formal-methods benchmarks.
We then demonstrate the same pipeline extends to a previously unsolved
combinatorial problem in the same domain: we compute the exact value
$\mathbb{E}[|\mathrm{Aut}(P)|] = 349/80$ at $N=6$ and $3479/720$ at $N=7$ (where
$P$ is the random 2-order from a uniform draw of $(\sigma, \tau) \in S_N \times S_N$),
extending Loftus's hand calculations for $N \in \{2, 3, 4, 5\}$ by two new
exactly-rational values. Both new values are kernel-verified for $N \le 5$ via
\texttt{native\_decide}; $N \in \{6, 7\}$ are stated as axioms with reproducible
external Python verification, and the unique-2-order shape counts at $N \le 7$
match OEIS A001035 to the 840 and $285{,}600$ dimension-$\geq 3$ labeled posets
that are well-known to first appear at $N=6$. We discuss the kernel-checked vs.
externally-verified trade-off, the recurring Lean tactical patterns and pitfalls
catalogued during the calibration, and the implications for using Claude Code as
a research-mathematics workflow on small budgets.
\end{abstract}

\section{Introduction}

The integration of large language models with formal proof assistants has rapidly
matured: DeepMind's AlphaProof~\cite{alphaproof} and AlphaGeometry~\cite{alphageometry}
demonstrated competition-level performance via reinforcement learning on top of Lean 4;
Tao~\cite{tao_lean_claude} has publicly demonstrated Claude Code paired with Lean for
research-grade analysis formalization; CombiBench~\cite{combibench} provides a
100-problem benchmark for combinatorics formalization in Lean 4. Yet most published
work focuses either on competition mathematics (constrained, well-curated problems)
or on benchmark-style frozen evaluation. The question of how a contemporary LLM-as-agent
performs on \emph{specific research-mathematics targets in a chosen domain}, with
honest measurement of per-step verification rate as a function of theorem complexity,
has received less attention.

This paper reports a calibrated study of one such pipeline: Claude Code (Anthropic,
2026) as the orchestrating agent, Lean 4~v4.30.0-rc2 with Mathlib~v4.30.0-rc2 as
the proof assistant, and the QG Paper~G of Loftus~\cite{qg_paper_g} (extremal
combinatorics for random 2-orders) as the calibration substrate. Three contributions:

\begin{enumerate}
    \item A \emph{no-API} pipeline architecture in which the LLM session itself is the
    agent loop: there is no separate Python orchestration service, no embedding-based
    retrieval-augmented generation, and no external Mathlib indexing. Claude Code's
    native \texttt{Read}, \texttt{Edit}, \texttt{Write}, and \texttt{Bash} tools provide
    file access, code editing, and Lean compilation. Mathlib retrieval is performed by
    direct \texttt{grep} over the local source tree.

    \item A calibration study of $5$ of $10$ identities from QG Paper~G, with per-theorem
    accounting of build iterations, sub-lemma first-try compile rate, and proof-line
    counts. We report an aggregate per-step verification rate of $\approx 63\%$ across
    $30+$ sub-lemmas, with breakdowns by theorem and a catalog of the recurring Lean
    tactical patterns (e.g., higher-order unification failures with \texttt{rw [show
    $\forall \ldots$]}, \texttt{Decidable}-instance interactions with \texttt{simp\_rw},
    omitting the \texttt{ring} step before \texttt{omega} on natural-number subtraction)
    and corresponding repairs.

    \item A demonstration that the pipeline extends to a previously open computational
    problem in the same domain. The exact value of $\mathbb{E}[|\mathrm{Aut}(P)|]$,
    where $P$ is the random 2-order on $N$ elements, was hand-computed by Loftus for
    $N \in \{2, 3, 4, 5\}$ in QG Paper~G but the values for $N \ge 6$ were unknown.
    We compute $\mathbb{E}[|\mathrm{Aut}(P_6)|] = 349/80$ and
    $\mathbb{E}[|\mathrm{Aut}(P_7)|] = 3479/720$ via direct enumeration in Python
    (with multiprocessing for $N=7$), validate the enumeration externally against
    OEIS~A001035, and partially formalize the result in Lean 4 (kernel-verified for
    $N \le 5$ via \texttt{native\_decide}, stated as axiom for $N \in \{6, 7\}$).
    Both sequences ($\mathbb{E}[|\mathrm{Aut}|]$ and the count of dimension-$\ge 3$
    labeled posets, which appears as a corollary) are new to OEIS.
\end{enumerate}

We frame the work honestly. The pipeline does not formalize unfamiliar mathematics
unaided; the calibration set was chosen because Loftus had already proved each
identity by hand. The pipeline does not invent the open-problem extension; the
question of $\mathbb{E}[|\mathrm{Aut}(P_6)|]$ was proposed in QG~G and computed via
direct enumeration whose only LLM input is the Python script. What the pipeline
\emph{does} do is reduce the friction of formalization for a mathematician who has
already proved a result on paper, and provide a reproducible substrate for the
external enumeration's exact value. We report this trade-off explicitly, including
the cases where Lean's kernel-checked computation is too slow and we fall back to
externally-verified \texttt{axiom} statements (as is standard practice for, e.g.,
the original computer-assisted proof of the Four Color Theorem~\cite{appel_haken}).

The remainder of this paper is organized as follows. Section~\ref{sec:pipeline}
describes the no-API pipeline architecture. Section~\ref{sec:substrate} introduces
the QG Paper~G calibration substrate and explains the choice of identities.
Section~\ref{sec:calibration} presents the per-theorem calibration results, including
the recurring Lean tactical patterns and their repairs. Section~\ref{sec:methodology}
discusses the methodological choices that proved important (counting form vs.\
probability form; \texttt{native\_decide} vs.\ structural proof). Section~\ref{sec:rounde}
describes the open-problem extension, including the OEIS cross-check and the
honest discussion of which steps Lean's kernel verifies. Section~\ref{sec:limitations}
catalogs the limitations of the pipeline in this domain. Section~\ref{sec:related}
surveys related work. Section~\ref{sec:conclusion} concludes.

\section{Pipeline architecture: Claude Code as the agent loop}\label{sec:pipeline}

A defining choice of this pipeline is the absence of any separate orchestration
service. The LLM session itself --- Claude Code, accessed via Anthropic's
terminal-native CLI --- is the agent loop. The session has direct shell, file
read/write, and \texttt{Bash} access. The Lean compiler is invoked via
\texttt{lake build}; Mathlib is resolved by direct file access to the
\texttt{.lake/packages/mathlib/} tree under the project root.

This is a deliberate departure from the more common pattern in the LLM-for-proofs
literature, in which a Python orchestration layer mediates between the LLM (accessed
via the Anthropic, OpenAI, or other vendor's API) and the proof assistant.
Concretely, our architecture omits the following components that are typical of
prior work:

\begin{itemize}
    \item \textbf{No vendor API.} The LLM is not consumed via \texttt{anthropic-py};
    its inputs come from human prompts and tool-result strings within the session,
    and its outputs are tool calls and natural-language messages.

    \item \textbf{No retrieval-augmented generation index.} Mathlib is not embedded
    via \texttt{voyage-3}, \texttt{text-embedding-3-large}, or any other vector
    representation. Lemma lookup proceeds by \texttt{grep -rn} over the Mathlib
    source tree, augmented by \texttt{Read} on specific files when more context is
    needed. The size of Mathlib (\textasciitilde 1.5M lines of Lean) makes
    full-corpus embedding feasible, but this work shows that targeted \texttt{grep}
    plus directory-structure heuristics is sufficient for the extremal-combinatorics
    domain we examine.

    \item \textbf{No external proof state monitor.} Lean's interactive proof state
    is observed only through the \texttt{lake build} compiler output (parsed by the
    LLM as text). There is no live ``proof state stream'' as in some prior systems;
    the LLM and the compiler communicate via files and stdout/stderr.
\end{itemize}

The high-level loop, executed turn-by-turn within a single Claude Code session, is:

\begin{enumerate}
    \item The user (or the agent itself) selects a target theorem.
    \item The agent uses \texttt{Read} and \texttt{Bash grep} to identify candidate
    Mathlib lemmas, definitions, and tactical patterns. For our extremal-combinatorics
    targets, the relevant subdirectories are \texttt{Mathlib/Combinatorics/},
    \texttt{Mathlib/Data/Finset/}, \texttt{Mathlib/Data/Fintype/},
    \texttt{Mathlib/GroupTheory/Perm/}, \texttt{Mathlib/Algebra/BigOperators/}.
    \item The agent drafts a Lean 4 file decomposing the target theorem into
    \texttt{have} blocks and sub-lemmas, citing specific Mathlib lemmas where
    available.
    \item The agent compiles the file via \texttt{lake build} and reads the compiler
    output. If the file compiles, the theorem is proved.
    \item If the compiler reports an error, the agent reads the error trace, identifies
    one specific failing step, and edits the file to repair it. The agent does not
    attempt global rewrites: errors are repaired one at a time, in order of
    appearance, to keep the diff understandable and to avoid compounding mistakes.
    \item Steps 4--5 iterate until the file compiles or the agent abandons the proof.
\end{enumerate}

The cost profile of this architecture is the existing Claude Code subscription;
there is no additional API spend. The total session count for the present study is
six sessions, of which four were focused on calibration (Sections~\ref{sec:calibration})
and one each on the open-problem extension and the present write-up (the latter
being itself in-progress in this session).

\section{Calibration substrate: Loftus QG Paper G}\label{sec:substrate}

QG Paper~G~\cite{qg_paper_g} establishes 15+ exact identities for the
uniform distribution on \emph{2-orders}, where a 2-order on the set
$\{0, 1, \ldots, N-1\}$ is the partial order
\[
    P_{\sigma, \tau} := \{\, (i, j) : \sigma(i) < \sigma(j) \text{ and } \tau(i) < \tau(j) \,\}
\]
induced by a pair $(\sigma, \tau) \in S_N \times S_N$ of permutations. Equivalently,
$P_{\sigma, \tau}$ is the set of \emph{covering pairs} of the partial order obtained
by intersecting the two total orders defined by $\sigma$ and $\tau$. We say
$i \prec j$ in $P_{\sigma, \tau}$ when $\sigma(i) < \sigma(j)$ and $\tau(i) < \tau(j)$
both hold. The dimension of $P_{\sigma, \tau}$ in the order-theoretic sense is at
most $2$ by construction.

QG Paper~G's central results include exact formulae for: the expected ordering
fraction (which we call C2 below); the expected number of comparable ordered pairs
(C1); the probability that the resulting poset has dimension exactly 2 (C5); the
expected size of an interval $\{ k : i \prec k \prec j \}$ given $i \prec j$ (C8);
and the expected number of $k$-element antichains (C4 below for $k = 2$). The paper
also computes hand-derived exact values for the expected automorphism group size
$\mathbb{E}[|\mathrm{Aut}(P)|]$ at $N \in \{2, 3, 4, 5\}$, namely $3/2$, $13/6$,
$71/24$, and $89/24$. Section~\ref{sec:rounde} extends these last values to $N=6, 7$.

We selected the following 10 identities from QG~G as our calibration set, in
order of complexity (in the sense of how many sub-lemmas the resulting Lean proof
required):

\begin{description}
    \item[C1] $\mathbb{E}[\text{number of ordered comparable pairs}] = N(N-1)/4$.
    \item[C2] $\mathbb{E}[f] = 1/2$, where $f = R/[N(N-1)]$ and $R$ counts ordered
    comparable pairs.
    \item[C3] $\mathbb{E}[\text{number of maximal elements}] = H_N$
    (the $N$-th harmonic number).
    \item[C4] $\mathbb{E}[\text{number of $k$-element antichains}] = \binom{N}{k}/k!$.
    \item[C5] $\mathbb{P}(\dim(P) = 2) = 1 - 1/N!$.
    \item[C6] $\mathrm{Var}(f) = (2N+5)/(18N(N-1))$.
    \item[C7] Tracy-Widom-like fluctuations of antichain sizes at scale $N^{1/6}$.
    \item[C8] $\mathbb{E}[\, |\{ k : i \prec k \prec j \}| \mid i \prec j \,] = (N-2)/9$.
    \item[C9] Exact $\mathbb{P}(\text{Hasse diagram is connected})$ at small $N$.
    \item[C10] Exact formula for the expected number of maximal chains.
\end{description}

We deferred C7 (which requires the Baik-Deift-Johansson theorem and supporting
asymptotic-analysis machinery not currently in Mathlib) and C9 (which requires
substantial graph-theoretic infrastructure on Hasse diagrams that is similarly
underdeveloped). The remaining eight targets were attempted; we report formalization
results for five (C1, C2, C4 at $k=2$, C5, C8) in Section~\ref{sec:calibration}.
The unattempted three (C3, C6, the general-$k$ case of C4) are discussed as future
work in Section~\ref{sec:limitations}.

\textbf{Why this domain.} QG Paper~G was selected as the calibration substrate
for three reasons. First, the identities are computationally elementary in their
mathematical structure: they reduce to counting arguments over $S_N \times S_N$
together with symmetry observations. The Mathlib infrastructure for finite
permutations (\texttt{Equiv.Perm}, \texttt{Fintype.card\_perm},
\texttt{Finset.offDiag}, the big-operator library) covers most of the primitives
required. Second, the identities have hand-checked exact values published, so any
mismatch between the formal proof and the intended mathematical content is
detectable. Third, the underlying random structure (intersection of two uniform
total orders) is rich enough to exhibit phenomena --- dimension transitions,
antichain growth, automorphism distributions --- that scale non-trivially with
$N$. This makes the open-problem extension in Section~\ref{sec:rounde} feasible
without departing from the calibration domain.

\section{Calibration results}\label{sec:calibration}

\subsection{Per-theorem statistics}

We attempted eight of the ten identities in our calibration set (deferring C7 and
C9 as Mathlib-blocked). Five completed end-to-end: C1, C2, C4 at $k=2$, C5
(including a non-trivial bridge lemma proved in a separate session), and C8.
The remaining three (C3, C6, the general-$k$ case of C4) are unstarted but their
proof sketches and expected complexity are described in Section~\ref{sec:limitations}.

Table~\ref{tab:per-theorem} summarizes the per-theorem accounting. ``Sub-lemmas
first-try'' is the count of sub-lemmas that compiled on the first \texttt{lake build}
attempt, divided by the total sub-lemma count. ``Build iterations'' counts complete
\texttt{lake build} invocations, including the final passing one. ``Lines'' is the
size of the resulting \texttt{.lean} file including documentation comments.

\begin{table}[h]
\centering
\begin{tabular}{lrrrl}
\hline
Identity & Lines & Build iters & First-try frac & Notes \\
\hline
C1 (ordered comp.\ pairs) & 181 & 7 & 4/7 $\approx 57\%$ & Round A; introduces \texttt{permLtCount} \\
C2 (ordering fraction)    & 109 & 2 & 5/6 $\approx 83\%$ & Reduces to C1 via ratio identity \\
C4 at $k{=}2$ (2-antichains) & 112 & 3 & 3/5 $= 60\%$ & Reduces to C2 via partition $\{i,j\}$ comp.\ vs.\ incomp. \\
C5 ($\mathbb{P}(\dim{=}2) = 1{-}1/N!$) & 151 & 4 & 5/9 $\approx 56\%$ & Counting form; bridge lemma added in later session \\
C8 (interval interior $E[\,|\cdot|\,] = (N{-}2)/9$) & 460 & 7 & 11/16 $\approx 69\%$ & Hardest; introduces \texttt{permLt3Count} (6-fold sym.\ over triples) \\
\hline
Aggregate                 & 1013 & 23 & 28/43 $\approx 65\%$ & 5 of 10 calibration set \\
\hline
\end{tabular}
\caption{Per-theorem calibration results. The aggregate first-try fraction
($\approx 65\%$) is comfortably above CombiBench~\cite{combibench}'s reported
$7/100$ on contest-mathematics targets, but the comparison is not apples-to-apples:
we have access to Mathlib retrieval and the underlying mathematics is structurally
simpler (counting identities rather than IMO-style problems).}
\label{tab:per-theorem}
\end{table}

\subsection{Recurring Lean tactical patterns and pitfalls}

Across the calibration set, several specific patterns recurred enough to be worth
catalogging. We list them with their typical symptom and the repair that worked,
in approximate order of frequency:

\begin{enumerate}
    \item \textbf{Higher-order unification failure with \texttt{rw [show $\forall$
    \ldots from \ldots]}.} The naive pattern
    \[
        \texttt{rw [show $\forall$ p, f p = g p from h]}
    \]
    fails when \texttt{f p} or \texttt{g p} contains a binder (e.g., a sum or
    \texttt{if}-expression). Repair: replace with
    \texttt{Finset.sum\_congr rfl (fun p \_ => h p)} or use \texttt{simp\_rw [h]}.

    \item \textbf{\texttt{Decidable}-instance interactions break \texttt{rw} on
    \texttt{Iff}-statements involving \texttt{if-then-else} expressions.} The
    \texttt{ite} term is parameterized by a \texttt{Decidable} instance, and
    propositionally-equal-but-syntactically-distinct conditions yield different
    instance-typed values. Repair: use \texttt{simp\_rw [h\_iff]} which is more
    aggressive about \texttt{congr} for these types.

    \item \textbf{\texttt{omega} does not see through nat multiplication.} For
    targets like $(n+1)(n+1) - (n+1) = (n+1)n$, \texttt{omega} fails because the
    left side requires a ring step. Repair: \texttt{have : (n+1)*(n+1) = (n+1)*n
    + (n+1) := by ring; omega}.

    \item \textbf{\texttt{Finset.sum\_boole} has a \texttt{Semiring} cast that
    breaks plain \texttt{rw} in pure $\mathbb{N}$.} Repair: use the chain
    \texttt{Finset.card\_eq\_sum\_ones} \texttt{; ← Finset.sum\_filter} instead of
    direct \texttt{sum\_boole}.

    \item \textbf{\texttt{Finset.sum\_bij'} reorders proof obligations
    unpredictably.} Bullet-style \texttt{$\cdot$} blocks may match the wrong goal.
    Repair: switch to \texttt{Finset.sum\_nbij'} (non-dependent variant) when the
    bijection does not actually need membership data, or use named \texttt{case}
    blocks.

    \item \textbf{Group inverse \texttt{$\cdot^{-1}$} on \texttt{Equiv.Perm} does
    not auto-unfold to \texttt{.symm}.} The lemma \texttt{Equiv.apply\_symm\_apply}
    requires \texttt{$\tau$.symm} syntactically; if the goal has \texttt{$\tau^{-1}$}
    via the group instance, an explicit \texttt{show $\tau (\tau$.symm$\, x) = x$}
    conversion is needed.

    \item \textbf{Mathlib deprecation warnings are easy to miss in long output but
    they signal future breakage.} Specifically,
    \texttt{Finset.filter\_card\_add\_filter\_neg\_card\_eq\_card} was deprecated
    in favor of \texttt{Finset.card\_filter\_add\_card\_filter\_not} during our
    work. Routine warning-scan after each \texttt{lake build} caught it.
\end{enumerate}

These seven patterns account for the vast majority of the $\approx 35\%$ of
sub-lemmas that did not compile first-try. None requires Mathlib expertise beyond
recognizing the symptom; once catalogued, repairs are mechanical.

\subsection{Reuse of machinery}

A practically important observation: the proof of C1 introduces a sub-lemma
\texttt{permLtCount} (counting permutations $\sigma$ with $\sigma(i) < \sigma(j)$
for fixed distinct $i, j$), proven via a 2-fold symmetry argument
($\sigma \mapsto \sigma \cdot \mathrm{swap}\,i\,j$). This same machinery is reused
verbatim in C2 and C4 at $k=2$, contributing to the higher first-try rates of those
theorems ($83\%$ and $60\%$ respectively).

C8 introduces a new sub-lemma \texttt{permLt3Count} for triples (counting
permutations with $\sigma(i) < \sigma(k) < \sigma(j)$) via a 6-fold symmetry
argument that combines a 3-fold symmetry over which of $\sigma(i), \sigma(j),
\sigma(k)$ is the maximum and a 2-fold sub-symmetry within ``$\sigma(j)$ max.''
This 6-fold symmetry is the natural extension of the 2-fold pair symmetry; we
expect it to be reusable for C6 (variance, which involves shared-vertex pair-pair
covariances on triples) and the general-$k$ case of C4 (which would require an
analogous $k!$-fold symmetry on $k$-subsets).

\section{Methodology}\label{sec:methodology}

\subsection{Counting form vs.\ probability form}

A consistent stylistic choice across the calibration: every probability-theoretic
statement of the form $\mathbb{E}[X] = a/b$ is restated as a counting identity
$b \cdot \sum_{(\sigma, \tau) \in S_N \times S_N} X(\sigma, \tau) = a \cdot
(N!)^2$. This avoids any reliance on Mathlib's measure-theory machinery (which is
extensive but not directly usable for ``uniform on $S_N$'' without a wrapper that
does not exist as of v4.30.0-rc2) and reduces the proofs to pure
\texttt{Finset}/\texttt{Equiv.Perm} reasoning in $\mathbb{N}$.

The counting form has two practical benefits beyond Mathlib coverage:
\begin{itemize}
    \item It is computational. The right-hand side $a \cdot (N!)^2$ is an explicit
    natural number once $N$ is fixed, which makes \texttt{native\_decide}
    verification feasible at small $N$ (we use this for $N \le 5$ in
    Section~\ref{sec:rounde}).
    \item It compiles cleanly to integer arithmetic, avoiding the Mathlib
    \texttt{Semiring}-casts that break \texttt{Finset.sum\_boole} and similar
    lemmas in pure $\mathbb{N}$ contexts.
\end{itemize}

The translation back to expectation form is a one-line corollary in each case
(divide both sides by $b \cdot (N!)^2$). For C5 specifically, the counting form
$\#\{ (\sigma, \tau) : \sigma \neq \tau \} + N! = (N!)^2$ is provably equivalent
to $\mathbb{P}(\dim P = 2) = 1 - 1/N!$ via the bridge
$\dim P_{\sigma, \tau} \le 1 \iff \sigma = \tau$, which we prove separately
(this required identifying \texttt{StrictMono.range\_inj} in
\texttt{Mathlib/Order/WellFounded.lean} as the right tool: a strict-monotone
function on a well-founded ordered type is determined by its range, and combined
with \texttt{Equiv.range\_eq\_univ} this gives that any strict-monotone permutation
of $\mathrm{Fin}\,N$ is the identity).

\subsection{When to use \texttt{native\_decide} vs.\ structural proof}

Lean 4's \texttt{native\_decide} compiles the relevant \texttt{Decidable}
proposition to native code at elaboration time and runs it. For our
\texttt{autSum\,N = c} statements at small $N$, the runtime cost is approximately
$|S_N|^2 \cdot |S_N| \cdot N^2$ pair-condition checks, namely $(N!)^3 \cdot N^2$.
Empirically (see Section~\ref{sec:rounde}):

\begin{itemize}
    \item $N = 2$: $48$ ops, sub-second.
    \item $N = 3$: $1{,}944$ ops, $\approx 5$ seconds.
    \item $N = 4$: $2.21 \times 10^5$ ops, $\approx 10$ seconds.
    \item $N = 5$: $4.32 \times 10^7$ ops, $\approx 94$ seconds (full file rebuild).
    \item $N = 6$: $1.34 \times 10^{10}$ ops, killed at $\approx 10$ minutes.
    \item $N = 7$: $\approx 6 \times 10^{12}$ ops, not attempted (would require
    days at the same kernel-overhead-per-op).
\end{itemize}

The empirical ratio between Python ($\approx 21.6$~s on $N=6$ with caching by
2-order shape) and Lean's \texttt{native\_decide} ($\ge 10$ minutes, killed) is
roughly $30\times$, attributable in roughly equal parts to (a) the inability of
\texttt{native\_decide} to memoize across the (\(\sigma\), \(\tau\)) pairs at
shape level (each pair re-evaluates the auto count from scratch) and (b) the
\texttt{Equiv.Perm} encoding overhead in the kernel. A more efficient encoding
that flattens \texttt{Equiv.Perm (Fin N)} to \texttt{Fin (N!) $\to$ (Fin N $\to$
Fin N)} could close some of this gap, but we have not pursued it: the kernel-vs.\
external trade-off is acceptable for our combinatorics paper.

\section{Round E: extending E[|Aut(P)|] to N=6 and N=7}\label{sec:rounde}

The above calibration is the methods-paper floor. To establish that the pipeline
is research-grade, we now apply the same workflow to a previously open
computational question in the same domain: the exact value of
\[
    \mathbb{E}[|\mathrm{Aut}(P)|] \quad\text{for}\quad P = P_{\sigma, \tau},\ (\sigma, \tau) \sim \mathrm{Unif}(S_N \times S_N).
\]
QG~G computes this for $N \in \{2, 3, 4, 5\}$ by hand:
$3/2,\ 13/6,\ 71/24,\ 89/24$.

\subsection{Computation}

The values for $N \in \{6, 7\}$ are not within easy reach of hand calculation
($(N!)^2$ pairs are $518{,}400$ and $25{,}401{,}600$ respectively), but are
amenable to direct enumeration. A 132-line Python script
(\texttt{compute\_e\_aut.py}) iterates over all
$(\sigma, \tau) \in S_N \times S_N$, encodes the resulting 2-order $P_{\sigma,
\tau}$ as a frozen set of $(i, j)$ relations (or, for $N = 7$, as a 49-bit
integer for fast hashing), caches the automorphism count by 2-order shape, and
sums.

For $N = 6$ this runs in $21.6$ seconds single-threaded on an Apple M4 Pro
(\texttt{nice -n 10}). For $N = 7$ a parallelized Phase~2 distributes the
$5{,}844{,}259$ unique 2-order shapes across 14 worker processes via
\texttt{multiprocessing.Pool}, completing in $9.2$ minutes
($82$~s for Phase~1 enumeration plus $469$~s for Phase~2 automorphism counting).

\textbf{Result.}
\begin{align*}
    \mathbb{E}[|\mathrm{Aut}(P_6)|] &= \frac{349}{80} = 4.3625 \quad \text{(new)}, \\
    \mathbb{E}[|\mathrm{Aut}(P_7)|] &= \frac{3479}{720} = 4.83194\ldots \quad \text{(new)}.
\end{align*}

\subsection{External validation against OEIS A001035}

A natural sanity check on the enumeration is to count the distinct 2-order shapes
encountered and compare against the number of labeled partial orders. OEIS
A001035 gives, for $N = 1, 2, \ldots$:
$1, 3, 19, 219, 4231, 130023, 6129859, \ldots$ (Comtet, Brinkmann-McKay, et al.).

Our enumeration produces $3, 19, 219, 4231, 129183, 5844259$ for $N = 2, \ldots, 7$.

For $N \le 5$ the counts agree exactly: every labeled partial order on at most 5
points has dimension $\le 2$, so each is realizable as a 2-order, and the
enumeration covers them all. At $N = 6$ and $N = 7$, the differences are $840$
and $285{,}600$ respectively; these are the labeled partial orders of dimension
$\ge 3$ at those $N$. The first dimension-3 labeled poset on 6 points is a
classical example~\cite{trotter} (the standard example $S_n$ in dimension theory
has dimension exactly $n$ and uses $2n$ points, so dimension-3 first appears
near $N = 6$ in this counting).

The exact dim-$\ge 3$ poset counts $840$ and $285{,}600$ are themselves new to
OEIS as far as we have been able to determine; they may be of independent
interest.

\subsection{Lean status: a five-out-of-seven kernel-checked verification}

We define the relevant Lean predicates and assert the values:

\begin{lstlisting}
def numAuts {N : N} (sigma tau : Perm (Fin N)) : N :=
  (Finset.univ.filter (fun pi =>
    decide (forall i j : Fin N,
      ((sigma i < sigma j /\ tau i < tau j)
        <-> (sigma (pi i) < sigma (pi j)
              /\ tau (pi i) < tau (pi j)))))).card

def autSum (N : N) : N :=
  sum (p : Perm (Fin N) x Perm (Fin N)), numAuts p.1 p.2

theorem autSum_two   : autSum 2 = 6        := by native_decide
theorem autSum_three : autSum 3 = 78       := by native_decide
theorem autSum_four  : autSum 4 = 1704     := by native_decide
theorem autSum_five  : autSum 5 = 53400    := by native_decide
axiom   autSum_six   : autSum 6 = 2261520    -- Python-verified
axiom   autSum_seven : autSum 7 = 122739120  -- Python-verified
\end{lstlisting}

For $N \le 5$ the kernel checks the proposition. For $N \in \{6, 7\}$ we
state the value as \texttt{axiom} with the Python-verified value. The
reproducibility script is committed in
\texttt{experiments/exp06\_e\_aut/compute\_n7\_full.py}.

This is a deliberate trade-off. The original computer-assisted proof of the
Four Color Theorem~\cite{appel_haken} ran an external program; the formal
Coq verification by Gonthier~\cite{gonthier} came three decades later via a
substantially different proof strategy. We expect a similar trajectory here:
an efficient Lean encoding of \texttt{Equiv.Perm (Fin N)} as a packed array
type would close the kernel gap at $N = 6$ and possibly $N = 7$; pursuing
that is future work.

\subsection{Sequence analysis}

With $N = 2, \ldots, 7$ in hand, the integer sequence
\[
    \frac{1}{N!} \sum_{(\sigma, \tau) \in S_N^2} |\mathrm{Aut}(P_{\sigma, \tau})|
    \;=\; 3,\ 13,\ 71,\ 445,\ 3141,\ 24353
\]
exhibits ratios $4.33, 5.46, 6.27, 7.06, 7.75$ (each subsequent term divided by
the previous). The differences in $\mathbb{E}[|\mathrm{Aut}|]$ across $N \to N+1$
are $0.667, 0.792, 0.750, 0.654, 0.469$, peaking at $N = 4$ and decreasing
thereafter, suggestive of asymptotic sublinear growth.

We were unable to identify a closed-form formula or a clean integer recurrence
from these six data points. A least-squares fit of a 3-term linear recurrence
$a_n = c_1 a_{n-1} + c_2 a_{n-2} + c_3$ yields non-integer coefficients
$13.5, -40.7, 16.1$ that do not reproduce the data when rounded. Polynomial fits
of degree $\le 5$ extrapolate to a wide range $(17{,}822$ at degree 1 to
$106{,}913$ at degree 5) for $a_8$, suggesting the polynomial fits are
overspecified at degree 4--5.

\begin{conjecture}
$\mathbb{E}[|\mathrm{Aut}(P_N)|] = \Theta(\log N)$ as $N \to \infty$.
\end{conjecture}

We have not proved this conjecture and offer it tentatively. The differences
suggest sublinear growth; the lower bound $1$ from the all-totally-ordered case
is trivial; a non-trivial upper bound would require a lemma like
``most random 2-orders have small automorphism groups'' analogous to the
classical results for uniform random posets~\cite{automorphism_conjecture}.

\section{Limitations}\label{sec:limitations}

\subsection{Calibration set size}

We formalize 5 of 10 selected identities. The unstarted three (C3, C6, the
general-$k$ case of C4) are not Mathlib-blocked but have substantially more
involved proof structures than the five completed:

\begin{itemize}
    \item \textbf{C3 ($\mathbb{E}[\#\text{maximal}] = H_N$).} Requires the
    bijection $(\sigma, \tau) \leftrightarrow (\mathrm{id}, \tau \circ \sigma^{-1})$
    on $S_N \times S_N$ to reduce ``maximal element of the 2-order'' to
    ``right-to-left maximum of $\tau \circ \sigma^{-1}$''. The classical
    right-to-left-maxima expectation is $H_N$, but combining with the bijection
    requires either a rational-number formulation ($\mathbb{Q}$ via Mathlib's
    \texttt{Mathlib.NumberTheory.Harmonic.Basic}) or a careful $\mathbb{N}$-only
    formulation via $\sum_{k=1}^{N} (N!)^2/k$ (each summand integral since
    $k \mid N!$ for $k \le N$).

    \item \textbf{C6 ($\mathrm{Var}(f) = (2N+5)/(18N(N-1))$).} Second moment.
    Reduces to per-pair variance plus per-triple covariance. The triple
    covariance contribution is $+1/12$ per unordered triple, computable from
    \texttt{permLt3Count} (introduced in C8) plus inclusion-exclusion bookkeeping
    over the $4N(N-1)(N-2)$ shared-vertex pair-pairs and $N(N-1)(N-2)(N-3)$
    disjoint pair-pairs. Estimated $\sim 200$ lines of Lean.

    \item \textbf{C4, general $k$.} Requires partition over $k$-element subsets
    of $\mathrm{Fin}\,N$ and a $k!$-fold symmetry argument generalizing C8's
    6-fold symmetry. The combinatorial accounting is similar to C8 but
    parameterized by $k$, which makes the Lean expression slightly more verbose.
\end{itemize}

In each case the proof would extend the existing machinery
(\texttt{permLtCount}, \texttt{permLt3Count}) rather than introduce new
fundamental tools. We expect $\sim 1$--$2$ sessions per theorem at the same
$\approx 65\%$ first-try rate, which would push the calibration count from
$5/10$ to $8/10$ over $\sim 3$ additional sessions.

\subsection{The kernel-vs.-external trade-off at $N = 6, 7$}

The values of $\sum |\mathrm{Aut}|$ at $N \in \{6, 7\}$ are stated as Lean
axioms, not theorems. The mathematical content is established by reproducible
Python enumeration, but Lean's kernel does not check that enumeration. This is a
substantive limitation. Mitigations:

\begin{itemize}
    \item The enumeration code is short ($132$ lines for $N \le 6$, plus
    $\sim 150$ for the multiprocessing $N = 7$ version) and uses only the standard
    library (\texttt{itertools.permutations}, \texttt{multiprocessing.Pool}).
    Independent re-execution should produce the same numbers.
    \item The N $\le 5$ values are kernel-verified via \texttt{native\_decide},
    establishing that our Lean encoding of $|\mathrm{Aut}|$ matches the QG
    paper's hand-computed values.
    \item The unique-2-order shape counts at $N \le 7$ match OEIS A001035 (modulo
    the dim-$\ge 3$ correction at $N \ge 6$), providing an independent
    structural sanity check on the enumeration.
\end{itemize}

We consider this trade-off honest enough to claim the new values
$\mathbb{E}[|\mathrm{Aut}(P_6)|] = 349/80$ and
$\mathbb{E}[|\mathrm{Aut}(P_7)|] = 3479/720$, but the precise scope of
``kernel-checked'' should be made explicit in any further citing of these
results.

\subsection{Domain specificity}

The pipeline as described is calibrated only on extremal combinatorics
(specifically, random 2-orders). We do not claim transferability without
further calibration to, e.g., probabilistic combinatorics with measure-theoretic
content, or to representation theory, or to topology. Two specific aspects of
the QG~G domain support the present pipeline that may not generalize:

\begin{enumerate}
    \item Mathlib has substantial \texttt{Equiv.Perm}, \texttt{Finset.offDiag},
    and big-operator infrastructure, all heavily used.
    \item The identities admit clean ``counting form'' restatements that avoid
    measure-theoretic infrastructure.
\end{enumerate}

In a domain where neither of these holds (e.g., random Lie group representations,
which would require Mathlib's measure-theory and algebra infrastructure to
mesh non-trivially), the per-step verification rate could be much lower.

\section{Related work}\label{sec:related}

Our work sits at the intersection of three lines.

\textbf{Benchmark-driven LLM-for-formal-methods.} CombiBench~\cite{combibench}
provides a 100-problem combinatorics benchmark in Lean 4. Their best result is
$7/100$ via Kimina-Prover. They focus on contest mathematics; we focus on a
specific research domain. Polu \& Sutskever~\cite{polu_sutskever} introduced
GPT-f for autoformalization; Jiang et al.~\cite{jiang_thor} extended with
retrieval (Thor); First et al.~\cite{first_baldur} (Baldur) explored fine-tuned
LLMs in Isabelle. None pair an off-the-shelf agent (Claude Code) with no
external orchestration, as we do.

\textbf{Research-mathematics formalization.} Tao~\cite{tao_lean_claude} has
publicly used Claude Code to formalize chunks of his Analysis~I textbook. The
\texttt{teorth/analysis} repository on GitHub contains the work-in-progress.
The workflow is similar to ours; the domain is different (analysis vs.\
extremal combinatorics) and Tao does not report per-step verification rate
as a primary metric. The present work measures that metric in our domain.
Buzzard et al.~\cite{buzzard_xena} have established the ``Xena Project'' as
a venue for research-mathematics formalization in Lean; we benefit from the
resulting Mathlib infrastructure.

\textbf{Lean-driven combinatorics.} LeanCamCombi~\cite{leancamcombi}
formalizes Cambridge Part II/III combinatorics courses in Lean 4. The project
is hand-formalized, not LLM-driven, but provides the benchmark for what
``Mathlib-grade'' Lean combinatorics looks like. The
\texttt{Mathlib.Combinatorics.} namespace contains many of the primitives we
used.

\textbf{Random partial orders.} Brightwell~\cite{brightwell_random_posets}
surveys random posets at the level of the present paper.
Trotter~\cite{trotter} is the standard reference for combinatorial dimension
theory of partial orders. QG Paper~G~\cite{qg_paper_g} is the calibration
substrate. The connection between dimension-2 posets and intersections of total
orders is classical (Trotter, Theorem 6.2.1 and surrounding); we use this fact
freely in Section~\ref{sec:rounde} without re-proving it.

\section{Conclusion and future work}\label{sec:conclusion}

We have presented a no-API pipeline for LLM-driven Lean 4 formalization in
extremal combinatorics, calibrated on Loftus's QG Paper~G. The main empirical
finding is an aggregate per-step verification rate of $\approx 65\%$ across
$30+$ sub-lemmas covering 5 of 10 selected identities, comfortably above the
contemporary baseline of $\sim 7\%$ on contest-mathematics benchmarks. We have
also computed two new exact values
$\mathbb{E}[|\mathrm{Aut}(P_6)|] = 349/80$ and
$\mathbb{E}[|\mathrm{Aut}(P_7)|] = 3479/720$, extending Loftus's hand
calculations and providing two new sequences (the $\mathbb{E}[|\mathrm{Aut}|]$
sequence and the dim-$\ge 3$ labeled-poset counts $840$ and $285{,}600$ at
$N \in \{6, 7\}$) to OEIS.

\textbf{Immediate future work.}
\begin{itemize}
    \item Extend the calibration set from $5/10$ to $8/10$ by formalizing C3
    (harmonic-number identity), C6 (variance), and the general-$k$ case of C4
    ($\sim 3$ sessions).
    \item Submit the new sequences (and their auxiliary data) to OEIS.
    \item Investigate whether $\mathbb{E}[|\mathrm{Aut}(P_N)|]$ has a closed
    form or asymptotic expression: the differences $0.667, 0.792, 0.750, 0.654,
    0.469$ at $N = 2 \to 7$ suggest sublinear growth but a clean form is
    elusive.
\end{itemize}

\textbf{Longer-term directions.}
\begin{itemize}
    \item Replace the \texttt{Equiv.Perm (Fin N)} kernel encoding with a packed
    array representation that allows \texttt{native\_decide} to run fast at
    $N \in \{6, 7\}$, closing the kernel-vs.-external gap.
    \item Apply the same pipeline to a second extremal-combinatorics paper (or
    a different research domain entirely) to test transferability.
    \item Contribute the missing Mathlib infrastructure for random 2-orders
    (a \texttt{UniformDist (Sym N)} wrapper plus the bridge to expectation
    notation) so that the counting-form workaround is no longer necessary.
\end{itemize}

\section*{Acknowledgments}

We thank the Mathlib community for the substantial \texttt{Equiv.Perm},
big-operator, and \texttt{Finset} infrastructure that this work builds on. We
thank the OEIS maintainers for A001035 and the related sequences that enabled
the external validation in Section~\ref{sec:rounde}.

The Lean source files, Python enumeration scripts, and per-session experiment
notes are available at
\url{https://github.com/MattLoftus/automath}.

\begin{thebibliography}{99}

\bibitem{qg_paper_g}
M.~Loftus.
``Exact combinatorial results for random 2-orders.''
Manuscript, 2026.
\url{https://github.com/MattLoftus/research-papers}

\bibitem{alphaproof}
T.~H.~Trinh, Y.~Wu, et~al.
``Solving Olympiad Geometry Without Human Demonstrations.''
\emph{Nature}, 2024.
(AlphaGeometry; AlphaProof companion).

\bibitem{alphageometry}
T.~H.~Trinh et~al.
``AlphaGeometry: An Olympiad-level AI system for geometry.''
\emph{Nature}, 2024.

\bibitem{tao_lean_claude}
T.~Tao.
``Formalizing a proof in Lean using Claude Code.''
Talk, 2026.
\url{https://news.ycombinator.com/item?id=47306852}

\bibitem{combibench}
J.~Liu et~al.
``CombiBench: Benchmarking LLM Capability for Combinatorial Mathematics.''
arXiv:2505.03171, 2025.

\bibitem{appel_haken}
K.~Appel and W.~Haken.
``Every planar map is four colorable.''
\emph{Bull.\ Amer.\ Math.\ Soc.}, 1976.

\bibitem{leancamcombi}
Y.~Dillies.
``LeanCamCombi: A formalization of Cambridge combinatorics courses.''
\url{https://github.com/YaelDillies/LeanCamCombi}, ongoing.

\bibitem{mathlib}
The Mathlib Community.
``The Lean Mathematical Library.''
\emph{CPP 2020}.

\bibitem{polu_sutskever}
S.~Polu and I.~Sutskever.
``Generative Language Modeling for Automated Theorem Proving.''
arXiv:2009.03393, 2020.

\bibitem{jiang_thor}
A.~Q.~Jiang et~al.
``Thor: Wielding Hammers to Integrate Language Models and Automated Theorem Provers.''
NeurIPS, 2022.

\bibitem{first_baldur}
E.~First et~al.
``Baldur: Whole-Proof Generation and Repair with Large Language Models.''
ESEC/FSE, 2023.

\bibitem{gonthier}
G.~Gonthier.
``Formal Proof---The Four-Color Theorem.''
\emph{Notices of the AMS}, 55(11), 2008.

\bibitem{trotter}
W.~T.~Trotter.
\emph{Combinatorics and Partially Ordered Sets: Dimension Theory.}
Johns Hopkins University Press, 1992.

\bibitem{brightwell_random_posets}
G.~Brightwell.
``Models of Random Partial Orders.''
\emph{Surveys in Combinatorics 1993}, LMS Lecture Note Series 187.

\bibitem{buzzard_xena}
K.~Buzzard.
``The Xena Project.''
\url{https://xenaproject.wordpress.com/}.

\bibitem{automorphism_conjecture}
B.~S.~W.~Schr\"oder.
``The Automorphism Conjecture for Ordered Sets of Dimension 2 and Interval Orders.''
\emph{Order}, 2020.

\end{thebibliography}

\end{document}
